\section{Splitting an odd number of roots}
\label{odd:section}

We now combine the geometry of Section~\ref{geom:section} with the divisor
algorithms.  The input is a square-free polynomial already known to split
into linear factors, but its roots are not given.  Computation over its
coordinate algebra treats all roots at once.  If a scalar zero test behaves
differently at two roots, it immediately produces a factor by
Lemma~\ref{ff:simultaneous}.  Proposition~\ref{geom:separation} supplies
the reason that a distinction must occur.  Our task here is to construct
its input classes and perform its label tests within the promised bound.

\begin{samepage}
\begin{proposition}[Odd-degree splitting]
\label{odd:split}
Let $p$ be prime, $L=\lceil\log_2p\rceil$, and
$B=20+(n+1)(L+1)$.  Suppose that $p>B^{200000}$ and that
$F\in\mathbb F_p[X]$ is monic, square-free, totally split, and of odd
degree $3\le N\le n$.  Choose an odd prime $q\mid N$.
Suppose that the following data are supplied explicitly:
\begin{itemize}
\item a field $K/\mathbb F_p$ of degree at most $q-1$;
\item a primitive $q$-th root of unity $\zeta\in K$;
\item an element $\omega_K\in K^*$ of order
      $q^{v_q(|K|-1)}$.
\end{itemize}
There is a deterministic algorithm, denoted by $\mathrm{OddSplit}$, that
returns a nontrivial proper monic divisor of $F$ over $\mathbb F_p$ in a
number of bit operations bounded by a fixed polynomial in $B$.
\end{proposition}
\end{samepage}

\subsection{The working field and the computational operations}

Put
\begin{equation}
             r=v_q(N),\qquad m=1+(q-1)r,
             \qquad Q_K=|K|,\quad s_K=v_q(Q_K-1).
\label{odd:parameters}
\end{equation}
Construct
\begin{equation}
             k=K[T]/(T^{q^r}-\omega_K),
             \qquad \omega=T\bmod(T^{q^r}-\omega_K).
\label{odd:working-field}
\end{equation}
This quotient is a field.  Indeed, every root of the binomial has exact
order $q^{s_K+r}$.  For every positive integer $b$, the elementary
lifting-the-exponent identity gives
\[
                  v_q(Q_K^b-1)=s_K+v_q(b).
\]
Thus the smallest extension of $K$ containing such a root has degree
$q^r$.  The binomial has that degree and is irreducible.
The same identity shows that $\omega$ generates the entire $q$-primary
subgroup of $k^*$.
Notice the bounds
\begin{equation}
                 q^r\le N,\qquad m\le q^r\le N,
                 \qquad [k:\mathbb F_p]\le N(q-1).
\label{odd:field-bounds}
\end{equation}
The middle inequality follows from
$q^r=(1+(q-1))^r\ge1+r(q-1)$.

Use the curve $C/K$ and its base extension to $k$ from
Section~\ref{geom:setup}.  Thus $Y^q=F(X)$,
$P_i=(x_i,0)$, $\lambda=1-\sigma$, and
$e=(\sigma-1)\infty_0$.  Work simultaneously over
\[
                    A=k[T_1]/F(T_1)\simeq\prod_{i=1}^N k.
\]
The element $T_1\bmod F$ provides the known ramification point in each
component computation.  The labels $i$ and values $x_i$ are used only in
the correctness proof.

Here are the precise divisor operations that the procedure uses.
The divisor arithmetic of Proposition~\ref{div:arithmetic} adds divisors,
applies $\sigma$, and reduces a degree-zero divisor to an equivalent one
of absolute degree at most $2g$.  Lemma~\ref{div:trace} sums a
simultaneously represented family across the components of $A$, without
finding those components.  Lemma~\ref{div:circuits} finds a defining
function for a principal divisor of the form $(1-\sigma)D-W$, even when
the infinity divisor $W$ has large integer coefficients.  It handles those
coefficients by addition and doubling circuits and retains the function
as a product circuit.  Valuations and leading values at $P_i$ can be
computed from this circuit.

The procedure $\mathrm{LambdaDivide}$ of
Proposition~\ref{div:divide} takes a reduced degree-zero divisor $D$
and a known ramification point $P_i$.  Its promise is that $[D]$ has a
$\lambda$-division represented by a divisor over $k$.  It returns a
reduced degree-zero divisor $D'$ over $k$ satisfying
\[
                              \lambda[D']=[D].
\]
All these procedures use field arithmetic and scalar zero tests, so they admit
the simultaneous execution just described.  At any test that distinguishes
components, $\mathrm{OddSplit}$ stops and returns the resulting factor.

\subsection{The procedure}

The algorithm has three stages: lift the ramification classes, sum the
lifts, and test labels.  We first give the complete instructions and then
prove that the promise for every call holds and that a factor must be
returned.

\begin{enumerate}
\item \textit{Lift each ramification class.}
In simultaneous execution over $A$, start with
$D=P_i-\infty_0$ in component $i$.
Apply $\mathrm{LambdaDivide}(D,P_i)$ successively $m-1$ times, each time
replacing $D$ by its output.  Write $d_{i,t}$ for the class after $t-1$
divisions, so $d_{i,1}=[P_i-\infty_0]$.

\item \textit{Sum the lifts and apply successive powers of $\lambda$.}
Use the simultaneous divisor-sum operation to obtain a divisor representing
\[
                              d_m=\sum_{i=1}^N d_{i,m}.
\]
Reduce it, obtaining $D_m$.
For $t=m-1,m-2,\ldots,1$, reduce
$(1-\sigma)D_{t+1}$ to obtain $D_t$.  Thus, writing $d_t=[D_t]$,
\begin{equation}
                           d_t=\lambda^{m-t}d_m.
\label{odd:downward-chain}
\end{equation}
Compute the integral infinity divisors
\begin{equation}
                           W_t=\lambda^{-(t-1)}Ne,
                           \qquad 1\le t\le m.
\label{odd:W}
\end{equation}
Here the inverse notation denotes the unique solution in the lattice
$\Lambda$, not division of divisor classes.  These solutions exist because
$q^r\mid N$ and $q$ is an associate of $\lambda^{q-1}$.
One may compute them successively by solving
$(1-\sigma)W_{t+1}=W_t$ on the cyclic infinity coordinates.

\item \textit{Test the labels.}
For $t=1,2,\ldots,m$, perform the following operations.
The relation $\lambda d_t=[W_t]$ will hold at every test that is reached.
Find a function $h_t\in k(C)^*$ satisfying
\begin{equation}
                      \operatorname{div}(h_t)
                                   =(1-\sigma)D_t-W_t.
\label{odd:h}
\end{equation}
It may be stored by the product-circuit method above.
Evaluate $h_t(P_i)$ and $v_{P_i}(D_t)$ simultaneously.
The value $h_t(P_i)^q$ is independent of $i$; extract it as a scalar
$c_t\in k^*$ from $A$.
Use $\mathrm{UnitRoot}$ from Lemma~\ref{ff:unitroot} to find
$\gamma_t\in k^*$ with $\gamma_t^q=c_t$.
For each component, compare $h_t(P_i)/\gamma_t$ with
$1,\zeta,\ldots,\zeta^{q-1}$ and form
\begin{equation}
                 l_{i,t}=v_{P_i}(D_t)
                        -\log_\zeta\bigl(h_t(P_i)/\gamma_t\bigr)
                            \pmod q.
\label{odd:labels}
\end{equation}
If these labels differ, return a nontrivial proper gcd with $F$
corresponding to one of their values.  If they all agree, proceed to the
next test.
\end{enumerate}

Every factor obtained by simultaneous execution lies in $\mathbb F_p[X]$.
Indeed, although the gcd may be computed over $k$, it is the monic product
of $X-x_i$ over a subset of roots $x_i\in\mathbb F_p$.
Thus no additional descent of the output factor is required.

\subsection{Correctness and cost}

\begin{proof}[Proof of Proposition~\ref{odd:split}]
Suppose that no earlier scalar test has returned a factor.
The elements $\zeta$ and $\omega$ supply the primary-group data required
by Theorem~\ref{norm:solve}.  In every division, its norm input is a
constant multiple of the function $j$ in Lemma~\ref{div:compression}.
Since the input divisor has absolute degree at most $2g$, the reduced
numerator and denominator of $j$ have degrees at most $2g$.
The norm solver's characteristic hypothesis therefore follows from
$(N+4g+1)^2<B^6<p$.
Corollary~\ref{geom:lifts}, with the $r,m,k$ of
\eqref{odd:parameters}--\eqref{odd:working-field}, proves the promise for
every invocation of $\mathrm{LambdaDivide}$ in the first stage.
This remains true for the particular divisions chosen by the algorithm:
every successive geometric division has a representative over $k$.
We therefore have
\[
                       \lambda^t d_{i,t}=[e]
                       \quad(1\le t\le m).
\]
The simultaneous sum in the second stage also satisfies its scalar-grid
requirement.  A reduced divisor has absolute degree at most $2g$, so its
compressed ideal bound has degree at most $2g$ and uses at most
$b=4qg+1$ generators.  Thus
\[
             N(b-1)\le 4Nqg
                  =2Nq(q-1)(N-2)<B^6<p,
\]
as required by Lemma~\ref{div:trace}.
Summing the equalities $\lambda^m d_{i,m}=[e]$ gives
\begin{equation}
                              \lambda^m d_m=[Ne].
\label{odd:sum-lifts}
\end{equation}

Equation~\eqref{odd:sum-lifts} is the hypothesis of
Proposition~\ref{geom:separation}.  The classes in
\eqref{odd:downward-chain}, their reduced representatives $D_t$, and
the lattice divisors in~\eqref{odd:W} are exactly its inputs.
That proposition proves that, at every reached test, the divisor in
\eqref{odd:h} is principal over $k$, the function is a unit at all
$P_i$, and its norm is a $q$-th power in $k^*$.
Lemma~\ref{div:circuits} computes $h_t$ and its values, and
Lemma~\ref{ff:unitroot} computes $\gamma_t$ with its promise thus
verified.  The logarithms in~\eqref{odd:labels} require only the stated
$q$ comparisons.
These are precisely the labels of Proposition~\ref{geom:separation},
which applies to all choices made by the deterministic routines.
It guarantees a nonconstant label vector by test $m$, unless an earlier
zero test has already returned a factor.  Therefore the procedure returns
a nontrivial proper factor.

For the running time, \eqref{odd:field-bounds} bounds the extension degree
and the number of divisions and tests by polynomials in $n$.
All divisors $D_t$ and all divisors retained between divisions are reduced,
with absolute degree at most $2g$.
The sole potentially large divisors are the $W_t$, and they have short
integer descriptions.  More explicitly, if
$W=\sum_j a_j\infty_j$, solving $(1-\sigma)W'=W$ amounts to taking
partial sums of the $a_j$ and subtracting their mean.  Whenever the
integral solution in $\Lambda$ exists, this increases the maximum
absolute coefficient by at most a factor $2q$.
Hence every coefficient of every $W_t$ has magnitude at most
$N(2q)^m$, and thus bit length $O(\log N+m\log q)$.
The principal-divisor computations use those coefficients through binary
circuits, as guaranteed by Lemma~\ref{div:circuits}.
All remaining loops and field operations have polynomial cost; the
explicit combined bit bound is derived in Section~\ref{sec:complexity}.
This proves the claimed uniform polynomial bound.
\end{proof}
